% GENERATED, edit the French unit or the translation table
\mxunit[9]{terminale-c05}{Combinatorics and counting}\label{ch:9}

\mxsubtitle{S22, S25, S30, S35}

\begin{keypoint}[Objectives]Choose a counting model; distinguish order and repetition; use binomial coefficients and Pascal's rule.\end{keypoint}

\section{Course notes}
\subsection{Adding or multiplying}
For pairwise disjoint finite sets $A_i$, $|\bigcup A_i|=\sum|A_i|$. If the cases overlap, the double counting must be corrected: $|A\cup B|=|A|+|B|-|A\cap B|$.
The number of elements of the Cartesian product $A\times B$ is $|A||B|$. For a sequence of choices, multiply the numbers of possibilities at each stage, provided these numbers are valid for every earlier choice considered.
\paragraph{Example.} A code made of two letters chosen from $26$, followed by three digits, with repetition allowed, can take $26^2\times10^3$ values. A three-digit number does not begin with $0$: there are $9\times10^2$ of them.
\subsection{Lists and subsets}
In a set of $n$ elements, there are $n^k$ ordered lists of length $k$ with repetition; $n(n-1)\cdots(n-k+1)=\frac{n!}{(n-k)!}$ lists of $k$ distinct elements, if $0\le k\le n$; $n!$ permutations of all the elements. We set $0!=1$.
\begin{keypoint}[Combinations]
A combination of $k$ elements from $n$ is a subset with $k$ elements; order does not matter.
\[\binom nk=\frac{n!}{k!(n-k)!},\quad\binom n0=\binom nn=1,\quad\binom n1=n,\quad\binom n2=\frac{n(n-1)}2.\]
\end{keypoint}
\paragraph{Why divide by $k!$?} Each subset of size $k$ gives exactly $k!$ distinct lists according to the order of its elements. The number of lists without repetition is therefore $k!\binom nk$.
\subsection{Identities and proofs}
\[\binom nk=\binom n{n-k},\qquad \binom nk+\binom n{k+1}=\binom{n+1}{k+1},\qquad\sum_{k=0}^n\binom nk=2^n.\]
Pascal's rule above is used for $0\le k<n$; the boundary formulas are read directly from the values $1$.
\paragraph{Symmetry.} Taking the complement sets up a correspondence between the subsets of size $k$ and those of size $n-k$.
\paragraph{Pascal, combinatorial proof.} In a set of $n+1$ elements, fix an element $a$. A subset of size $k+1$ either contains $a$ (choose $k$ other elements from $n$) or does not contain it (choose $k+1$ elements from $n$). The two cases are disjoint and cover all the possibilities.
\paragraph{Pascal, algebraic proof.} For $0\le k<n$, put over a common denominator:
\[\frac{n!}{k!(n-k)!}+\frac{n!}{(k+1)!(n-k-1)!}
=\frac{n!((k+1)+(n-k))}{(k+1)!(n-k)!}
=\frac{(n+1)!}{(k+1)!(n-k)!}.\]
The boundary cases are dealt with directly.
\paragraph{Number of subsets.} For each of the $n$ elements, choosing ``present'' or ``absent'' gives $2^n$ subsets. Sorting them by size, we count $\sum_{k=0}^n\binom nk$ of them: the two counts are equal.
\paragraph{Extension mentioned in the scheme of work.} The expansion $(a+b)^n=\sum_{k=0}^n\binom nk a^{n-k}b^k$ can be understood by choosing, in each of the $n$ factors, the term $a$ or $b$. This formula is not used as a prerequisite for the tests in this booklet.


\mxnewpage
\section{Research activity}\textbf{Order and repetition in counting --- 30 min.}\par Six distinct students volunteer to organise a debate.
\begin{enumerate}\item Two students are chosen to form a pair with no roles assigned. How many pairs can be formed?
\item A chair and a secretary are appointed, who must be two different students. How many assignments of the two roles are possible?
\item Three successive speeches are planned. The same student may speak several times. How many programmes of three speeches are possible?
\item Start again with only four students, denoted $A$, $B$, $C$ and $D$. Write out all the pairs and all the assignments of the two roles. Then explain why there are $64$ programmes of three speeches, without listing them.
\item For each of the three ways of choosing, state whether order matters and whether the same person may be chosen several times.\end{enumerate}

\mxnewpage
\section{Progressive exercises}

\mxexerciseitem{1}{Practice}A code is a sequence of four digits chosen from $0,1,\ldots,9$. Repetitions and the digit $0$ in first position are allowed. How many codes contain at least one $0$?

\mxexerciseitem{2}{Practice}Ten athletes take part in a race with no ties. How many different podiums, with a first, a second and a third place, can be obtained?

\mxexerciseitem{3}{Practice}From seven distinct people, a group of three people is formed containing one person designated in advance. How many groups can be formed?

\mxexerciseitem{4}{Practice}On a grid, we start from $\pt{0;0}$ and reach $\pt{4;3}$ using exactly four moves of one unit to the right and three upwards. How many different paths can be followed?

\mxexerciseitem{5}{Practice}A class has eight girls and seven boys. How many groups of four students containing exactly two girls can be formed?

\section{Textbook or past-paper work}\textbf{Barbazo, p. 294, exercise 37.} At home after S35: do both questions. State whether order matters and whether a card may be chosen several times before using a binomial coefficient.

\mxnewpage
\section{Homework}\textbf{Homework 9 --- Forming a committee --- 45 min, after S35.}\par A team consists of five teachers and seven students, all distinct. A committee of four people is formed.
\begin{enumerate}\item How many different committees can be formed?
\item How many committees contain exactly two students?
\item How many committees contain at least one teacher and at least one student?
\item Among the members of a committee from question 2, a teacher is appointed chair and a student secretary. Count the choices of the committee together with these two appointments.
\item By counting the committees according to their number of teachers, prove that \[\sum_{k=0}^4\binom5k\binom7{4-k}=\binom{12}4.\]\end{enumerate}

\ifmxteacher\mxnewpage\section{Teaching notes and solutions}\begin{teachernote}[Teacher notes]The common test in S28 comes before the combinations of S30: include only principles and ordered lists. Test 3 may involve combinations, but not Pascal's rule taught in S35. Have students state the object being counted before any formula. For weaker students, use cards bearing names; does permuting the cards change the object? Newton's binomial theorem, cited in the scheme of work, is an extension: keep it explicitly optional.\end{teachernote}

\subsection{Solutions to the activity}\begin{enumerate}
\item A pair is a subset of two students from six: $\binom{6}{2}=\frac{6\times5}{2}=15$.
\item The chair can be chosen in six ways, then the secretary from the other five students: $6\times5=30$ assignments.
\item Each speech can be given to any of the six students, including one already chosen. The three successive choices give $6^3=216$ sequences.
\item With four cards, the pairs are $AB,AC,AD,BC,BD,CD$: six choices. The twelve pairs of roles are:\[\begin{array}{lll}(A,B)&(A,C)&(A,D)\\(B,A)&(B,C)&(B,D)\\(C,A)&(C,B)&(C,D)\\(D,A)&(D,B)&(D,C)\end{array}\] There are $4^3=64$ sequences of three speeches, since each position has four choices.
\item For a pair, order does not matter and a student is not repeated. For the roles, order matters, with no repetition. For the speeches, order matters and repetition is allowed.
\end{enumerate}

\paragraph{E1.}There are $10^4$ codes in total. Those containing no $0$ have nine choices at each of the four positions, that is $9^4$ codes. By complement, the number required is $10^4-9^4=3439$.

\paragraph{E2.}There are ten choices for gold, then nine for silver and eight for bronze. The multiplication principle gives $10\times9\times8=720$ podiums. Changing the order changes the podium.

\paragraph{E3.}The imposed person is already chosen. It remains to choose a subset of two people from the other six: $\binom{6}{2}=\frac{6\times5}{2}=15$.

\paragraph{E4.}A path is determined by the positions of its four steps to the right among seven positions. The other three steps are then forced. We therefore count $\binom74=35$ paths.

\paragraph{E5.}We must choose two girls from eight, then two boys from seven. Each pair of choices gives a single group: $\binom82\binom72=28\times21=588$.

\mxneed{30}

\subsection{Solutions to Homework 9}\begin{enumerate}
\item Order does not matter: there are $\binom{12}{4}=\frac{12\times11\times10\times9}{24}=495$ committees.
\item We choose two students from seven and two teachers from five: $\binom72\binom52=21\times10=210$ committees.
\item From the $495$ committees, remove the $\binom54=5$ committees made up only of teachers and the $\binom74=35$ committees made up only of students. These categories are disjoint: $495-5-35=455$.
\item Each of the $210$ committees has two choices for the teacher chair and two choices for the student secretary: $210\times2\times2=840$ arrangements.
\item Sort the committees according to the number $k$ of teachers, with $0\le k\le4$. Category $k$ contains $\binom5k\binom7{4-k}$ committees. The five categories are disjoint and cover all the committees. By addition,
\[\sum_{k=0}^4\binom5k\binom7{4-k}=\binom{12}4=495.\]
\end{enumerate}

\subsection{Solutions to the textbook or past-paper work}A hand contains five distinct cards; order does not matter. There are therefore $\binom{52}5=\num{2598960}$ hands. For a hand containing the ace of spades, this card is already chosen; we choose the other four from the $51$ remaining cards. There are $\binom{51}4=\num{249900}$ hands. This result is not multiplied by $5$: the hand has no first or last position.

\fi

